The impact of the \gls{DNP} groups on transients is evaluated using Moltres with fluid-flow-velocity field coupled multi-group diffusion kinetics, as discussed in Section \ref{sec:mgd-kinetics}.
This approach provides high-fidelity results which couple the thermal-hydraulics and neutronics, a necessary condition for \gls{MSR} modeling.
Because modeling every \gls{DNP} in space and time would be prohibitively expensive, they are instead modeled using the \gls{DNP} group parameters.
This means that static \gls{DNP} group parameters which do not account for \gls{MSR} phenomena can be directly compared with the \gls{DNP} group parameters which do account for these phenomena.

The results of interest from this analysis will be the reactivity and power responses due to various initiating events.
The initiating event and transients of interest also vary, as there are different reactor responses that are of interest.
The events of interest are dictated by the perspective from which the reactor response is studied.

Specifically, there are four different perspectives of interest here: validation, safety, security, and safeguards.
Within each of these perspectives, there is a unique combination of motivation, proposed initiating event, and investigated response.
These are discussed specifically in each respective section.

\subsection{Validation}
The motivation behind validation is to confirm that the proposed methodology and results match experimental results.
This ensures that the demonstrated \gls{DNP} group parameters are useful for actual \glspl{MSR}.
Discrepancies between the experimental results and the model offer insight into potential weaknesses of the approach.
If there are no discrepancies, than it is reasonable to have more confidence in the results for problems similar to those which have been validated.

The validation of the \gls{DNP} group parameters has a multi-stage approach.
The first stage is validation of the irradiation by comparing the irradiated sample's generated delayed neutron count rate to that of the Keepin et al. experiment from 1957 \cite{keepin1957delayed}.
If the delayed neutron count rates do not match, then this means there is an issue in the data being used or the aggregation of the individual \glspl{DNP} into the count rate.
If the delayed neutron count rate is accurately predicted, then the delayed neutron count rate formulation can be accepted as valid.

The second stage of the validation is for the generated \gls{DNP} group parameters.
This stage uses the \gls{MSRE} pump start-up and coast-down experimental data to see if the group parameters accurately model the reactivity changes from the \glspl{DNP} \cite{prince1968zero}.
This experiment has been previously studied using Moltres and Quasimolto, so the results can be compared against them as well \cite{reynolds2023analysis, park_advancements_2025}.
However, those works did not include the inlet and outlet plena, which led to results which did not match the experimental data.
A model which includes the plena should give results closer to the experimental data, particularly when combined with the updated \gls{DNP} group parameters.

\subsection{Safety}
The motivation behind studying safety is to protect the public and meet regulatory requirements.
Safety is discussed in detail in Section \ref{sec:safety}, though the general takeaway is many possible safety-related transients exist.
The \gls{DNP} group parameters are important in these transients, as they affect the reactor period, power, and temperature.

The transients investigated are \gls{ULOHS}, \gls{ULOF}, \gls{UTOP}, and \gls{FSOC}.
For each of these transients, the reactor period, power response, maximum temperature, minimum temperature, and reactivity response are all recorded.
The reactor period is important, as it represents the controllability of the reactor.
The power response shows how changing the \gls{DNP} group parameters affects other responses, such as the temperature.
It is important to make sure differences in the \gls{DNP} parameters don't lead to temperature changes or oscillations which result in unsafe reactor conditions.

\subsection{Security}
Security is discussed in detail in Section \ref{sec:security}, but the general takeaway for this study is that bad actors may attempt to sabotage the reactor.
This ties directly into the safety design study, as the sabotage would be an attempt to create unsafe reactor conditions.
It is important, then, that the \gls{MSR} of interest does not have unsafe operations after a sabotage.
Particularly, if static \gls{DNP} group parameters predict safe operation while the updated \gls{DNP} group parameters predict unsafe operation, then this would be important to correct in the reactor design.

Because the safety study investigates most of the transients which would be of interest for sabotage, additional considerations can be included in the transients instead.
Because these transients are initiated by sabotage, it becomes more reasonable to then consider if multiple safety-related transients occur simultaneously. 
With four unique transients, this yields six unique combinations of double transients, three unique combinations of triple transients, and one combination of all four transients, for a total of ten security-related sabotage-initiated transients.

\subsection{Safeguards}
Safeguards is discussed in Section \ref{sec:safeguards}, but the main takeaway is that it is important to safeguard the \gls{MSR} of interest from having \gls{SNM} diverted.
One aspect of safeguarding the reactor is detection, which can take forms such as neutron detectors and surveillance systems \cite{zubair2024nuclear}.
The updated \gls{DNP} group parameters used in this work affect the effective delayed neutron fraction, and thus the reactivity response.
This means means detecting \gls{SNM} diversion by studying reactivity responses using static \gls{DNP} group parameters predicts less accurate responses than the updated \gls{DNP} group parameters.
If this difference is sufficiently large, it may be possible that failure to use these updated parameters would reduce the detection probability of a diversion scenario.

To investigate this, diversion scenarios are studied with a single significant quantity of material removed over the course of one year, simulated using OpenMC.
Then, the \gls{DNP} group parameters are generated for the fuel composition at that time, weighted by the volumetric fission rate fraction of each fissile nuclide, $F$, similar to Equation \eqref{eq:lami-multinuc}.
For example, the total delayed neutron yield would be given as:
\begin{equation}
    \bar{\nu}_{d} = \left(\sum_F \int \Sigma_f^{(F)} \phi \;dE d\vec{r} \right)^{-1}  \sum_F \left( \bar{\nu}^{(F)}_{d} \int \Sigma_f^{(F)} \phi \;dE d\vec{r} \right),
\end{equation}
where,
\begin{align}
    \bar{\nu}^{(F)}_{d} &= \text{ the total delayed neutron yield for nuclide $F$ [-]}\\
    \Sigma_f^{(F)} &= \text{ the macroscopic fission cross section for nuclide $F$ $[cm^{-1}]$}.
\end{align}
These weighted group parameters are generated using the static and the updated values.
Then, they can be used in simple step reactivity insertion transients, where the following four sets of data are compared: static group parameters without diversion, static group parameters with diversion, updated group parameters without diversion, and updated group parameters with diversion.
If the difference between the static group parameters with and without diversion is smaller or similar to the difference between the static and updated group parameters, then this means detection of that diversion scenario requires models to account for the \gls{MSR} physical phenomena when generating group parameters.
